\begingroup
\def\xmin{0} \def\xmax{0.055} \def\ymin{0} \def\ymax{3}
\def\maxxa{0.0199} \def\maxya{1.07} \def\maxxb{0.0483} \def\maxyb{2.65} % CHANGE: read off your graph
\def\minxa{0.0199} \def\minya{1.13} \def\minxb{0.0483} \def\minyb{2.55} % CHANGE: read off your graph
\pgfmathsetmacro{\gmax}{(\maxyb-\maxya)/(\maxxb-\maxxa)}
\pgfmathsetmacro{\gmin}{(\minyb-\minya)/(\minxb-\minxa)}
\pgfmathsetmacro{\gmean}{(\gmax+\gmin)/2}
\pgfmathsetmacro{\gunc}{abs(\gmax-\gmin)/2}

\section{Practical 1: Boyle's law}

\subsection{Aim}
To Investigate Boyle's Law
\subsection{Variables}
\begin{description}
  \item[Independent:] Volume(Length)
  \item[Dependent:] Pressure
  \item[Control:] Mass, Temperature
\end{description}

\subsection{Method}
\begin{enumerate}
    \item Given measuring the cross sectional area of the tube is rather difficult; the area stays constant throughout, therefore the length of the fluid in the tube is proportional to the volume of the fluid in the capillary tube.
    \item Set up the apparatus shown in the diagram, ensuring the tube is empty
    \item Drawing in air each time to the tube (slowly), record the height ($L_n$ /cm), until the total volume of the tube is filled up
    \item Repeat the experiment removing the water in the tube and record the mean volume for each mass ($\overline{L_n}$ /cm)
    \item Plot a graph of $P$ / \unit{\pascal} (y-axis) against $\frac{1}{L}$ / \unit{\per\centi\metre} (x-axis), drawing a best-fit line, and deduce from the straight line through the origin that $P \propto \frac{1}{L}$.
\end{enumerate}

\subsection{Risk assessment}
\begin{table}[H]\centering
\begin{tabularx}{\textwidth}{@{}l X X@{}}
\toprule
Risk & Hazard & Mitigation(s) \\
\midrule
\bottomrule
\end{tabularx}
\end{table}

\subsection{Results}

\begin{table}[H]\centering
\pgfplotstabletypeset[
  col sep=comma,
  columns={L,invL,P},
  columns/L/.style={column name={$L$ / \unit{\centi\metre}}, fixed, fixed zerofill, precision=1},
  columns/invL/.style={column name={$1/L$ / \unit{\per\centi\metre}}, fixed, fixed zerofill, precision=4},
  columns/P/.style={column name={$P$ / \unit{\pascal}}, fixed, fixed zerofill, precision=2},
  every head row/.style={before row=\toprule, after row=\midrule},
  every last row/.style={after row=\bottomrule}
]{data-boyles.csv}
\caption{Results table.}
\end{table}

\begin{figure}[H]\centering
\begin{tikzpicture}
\begin{axis}[
  width=0.85\textwidth, height=8cm,
  xlabel={$1/L$ / \unit{\per\centi\metre}}, ylabel={$P$ / \unit{\pascal}},
  xmin=\xmin, xmax=\xmax, ymin=\ymin, ymax=\ymax,
  grid=both, minor tick num=1,
  legend pos=north west,
]

\addplot[only marks, mark=x] table[col sep=comma, x=invL, y=P]{data-boyles.csv};
\addlegendentry{Data}

\addplot[blue, no marks] table[col sep=comma, x=invL,
  y={create col/linear regression={y=P}}]{data-boyles.csv};
\addlegendentry{Best fit}

\addplot[red, dashed, no marks] coordinates {(\maxxa,\maxya) (\maxxb,\maxyb)};
\addlegendentry{Max gradient}
\addplot[green!50!black, dashed, no marks] coordinates {(\minxa,\minya) (\minxb,\minyb)};
\addlegendentry{Min gradient}
\end{axis}
\end{tikzpicture}
\caption{Graph of $P$ against $\frac{1}{L}$ with best-fit line and worst-fit lines.}
\end{figure}

\subsection{Analysis}
Gradient of maximum line: $\num[round-mode=places,round-precision=2]{\gmax}$\\
Gradient of minimum line: $\num[round-mode=places,round-precision=2]{\gmin}$\\
Mean gradient: $\num[round-mode=places,round-precision=2]{\gmean}$\\
Uncertainty in gradient: $\pm\,\num[round-mode=places,round-precision=2]{\gunc}$

\begin{equation}
  P = k \cdot \frac{1}{L} \quad\Rightarrow\quad k = \text{constant}
\end{equation}

\subsection{Uncertainties}
\subsection{Evaluation}
\endgroup